\subsection{\textit{Sprint} 0}

Para a \textit{sprint} inicial, as atividades previstas envolviam o começo dos
repositórios de \textit{back-end} e \textit{front-end}, além da criação dos \textit{mockups} do Figma, da
modelagem inicial do banco de dados e, principalmente, de um intensivo de estudos
das tecnologias a serem usadas. Da minha parte, como AGES II, fiquei encarregado
de participar da criação dos \textit{mockups}, iniciar a modelagem do banco de dados e
estudar as tecnologias Python, FastAPI, PostgreSQL e Docker \cite{docker-docs}. Também foi prevista
uma \textit{spike} com o objetivo de definir quantas informações o grupo conseguiria
extrair a partir dos dados das empresas e como esses dados poderiam ser
extraídos.

A equipe conseguiu realizar muito bem as tarefas previstas e, da minha parte,
consegui superar minhas expectativas. Pesquisei \ac{api}s possíveis de extração,
principalmente nos períodos de aula. Além disso, participei ativamente da criação
dos \textit{mockups} no Figma e dediquei tempo ao estudo das tecnologias Python,
FastAPI e PostgreSQL.

O time como um todo trabalhou muito bem, com sincronia e boa comunicação. O
único possível problema relatado foi técnico: até o momento, ainda não havia
certeza se o projeto seria possível, visto que dependíamos de \ac{api}s externas de
redes sociais. Apesar disso, a equipe estava confiante quanto às chances de
conseguir coletar os dados necessários.

Acredito que, nesta primeira \textit{sprint}, tive muito a aprender com meus companheiros
de equipe, AGES IV e com o professor Daniel. Como gestores do projeto e da
equipe, eles realizaram um bom trabalho em reunir uma equipe que rapidamente
criou entrosamento e resultou em um ambiente de trabalho agradável. Por meio de
dinâmicas e votações, eles conseguiram unir o time inteiro.

Para a próxima \textit{sprint}, definimos como prioridade finalizar o banco de dados,
deixar os \textit{mockups} definitivos, definir os métodos de obtenção de dados e começar
a escrever código.
